\lyseis{Α1.1}{Δυνάμεις φυσικών αριθμών}
\sol{1} α) $5^4$ β) $7^3$ γ) $10^6$ δ) $x^4$ ε) $\alpha^2\cdot\beta^3$.
\sol{2} α) $32$ β) $81$ γ) $216$ δ) $100.000$ ε) $1$ στ) $0$ ζ) $225$.
\sol{3} Στο $6^7$ βάση $6$, εκθέτης $7$. Στο $15^2$ βάση $15$, εκθέτης $2$ (τετράγωνο του $15$). Στο $1^4$ βάση $1$, εκθέτης $4$. Στο $10^9$ βάση $10$, εκθέτης $9$. Κύβος δεν υπάρχει· «τετράγωνο» είναι μόνο η $15^2$.
\sol{4} $121,\ 144,\ 169,\ 196,\ 225$.
\sol{5} α) Λάθος ($16$). β) Λάθος ($8\neq 9$). γ) Σωστό. δ) Λάθος ($1^{20}=1$). ε) Σωστό. στ) Λάθος ($25\neq 10$).
\sol{6} α) $8<9$, άρα $<$. β) $64<81$, άρα $<$. γ) $100<1.024$, άρα $<$. δ) $1=1$, άρα $=$.
\sol{7} α) $20-12=8$. β) $17\cdot 4=68$. γ) $6\cdot 2=12$. δ) $36:12=3$. ε) $4+9=13$. στ) $7^2=49$.
\sol{8} $10\cdot 10\cdot 10\cdot 10=10^4=10.000$ κωδικοί. Με $6$ ψηφία: $10^6=1.000.000$ κωδικοί.
\sol{9} $3\cdot 3\cdot 3\cdot 3=3^4=81$ φυτά.
\sol{10} $12\cdot 12=12^2=144$ πλακάκια.
\sol{11} $4^3=4\cdot 4\cdot 4=64$ κυβικά δεκατόμετρα, δηλαδή $64$ λίτρα.
\sol{12} Τα ρέστα είναι $20-(4\cdot 2+3\cdot 3)=20-(8+9)=20-17=3$ €.
\sol{13} Το $2^5$ είναι $2\cdot 2\cdot 2\cdot 2\cdot 2=32$ και όχι $2\cdot 5=10$. Το $4^3=4\cdot 4\cdot 4=64$ και όχι $4\cdot 3=12$. Ο Νίκος πολλαπλασίασε βάση και εκθέτη, ενώ ο εκθέτης δείχνει πόσες φορές γράφεται η βάση.
\sol{14} α) $(2+3)\cdot 4=20$. β) $(12-4):2=4$. γ) $(2\cdot 3)^2=36$.
\sol{15} $10^{12}=1.000.000.000.000$ (το $1$ και $12$ μηδενικά). $2\cdot 10^5=200.000$ και $35\cdot 10^3=35.000$.
\sol{16} Τα τελευταία ψηφία των $2^1,\dots,2^8$ είναι $2,4,8,6,2,4,8,6$: επαναλαμβάνονται ανά τέσσερα. Ο εκθέτης $20$ είναι πολλαπλάσιο του $4$, όπως οι εκθέτες $4$ και $8$ που δίνουν $6$, άρα το $2^{20}$ τελειώνει σε $6$.
\sol{17} $5.082=5\cdot 10^3+0\cdot 10^2+8\cdot 10+2$ και $40.730=4\cdot 10^4+0\cdot 10^3+7\cdot 10^2+3\cdot 10+0$. Με ψηφία: $6.407$ και $200.080$.
\sol{18} $\alpha=8-(9-7)=8-2=6$. $\beta=4\cdot(9-8)+(4-4)^3=4\cdot 1+0=4$. $\mathrm{Α}=(6+4)^2-2\cdot 6\cdot 4=100-48=52$.
